\begin{table}[H]
    \centering
    \footnotesize
    \caption{Concordância entre LIME e SHAP nas três escolas discutidas no texto}
    \begin{tabular}{L{1.9cm}L{2.6cm}C{1.5cm}C{1.6cm}C{1.6cm}C{1.5cm}C{1.4cm}}
        \toprule
        \textbf{Trilha} & \textbf{Caso} & \textbf{Fidelidade (R\textsuperscript{2})} & \textbf{Estabilidade (Jaccard)} & \textbf{Top-10 comum (Jaccard)} & \textbf{Spearman} & \textbf{Mesmo sinal (\%)} \\
        \midrule
        Proporção & típica & 0,45 & 0,76 & 0,05 & 0,48 & 100 \\
        Proporção & acima do esperado & 0,73 & 0,87 & 0,05 & 0,52 & 100 \\
        Proporção & abaixo do esperado & 0,74 & 0,87 & 0,11 & 0,59 & 100 \\
        Proficiência & típica & 0,60 & 0,79 & 0,05 & 0,14 & 100 \\
        Proficiência & acima do esperado & 0,46 & 0,74 & 0,18 & 0,60 & 100 \\
        Proficiência & abaixo do esperado & 0,59 & 0,79 & 0,05 & 0,60 & 100 \\
        \bottomrule
    \end{tabular}
    \vspace{0.1cm}
    \begin{flushleft}
    {\footnotesize Fonte: elaborada pelo autor, com base nos microdados do Saeb 2023 (Inep). Fidelidade: R\textsuperscript{2} médio do modelo local do LIME. Estabilidade: Jaccard médio das 10 variáveis principais entre repetições do LIME. Top-10 comum: Jaccard entre as 10 principais do LIME e do SHAP. Spearman: correlação das importâncias. Mesmo sinal: variáveis comuns com contribuição de mesmo sinal.}
    \end{flushleft}
    \label{tab:lime}
\end{table}
